\subsection{Nhánh trong Git: con trỏ nhẹ và thao tác cơ bản}

Trong Git, nhánh (\textit{branch}) thực chất chỉ là một con trỏ nhẹ trỏ tới một commit cụ thể, thay vì một bản sao toàn bộ mã nguồn như nhiều hệ thống cũ \cite{progit}. Mỗi commit lưu con trỏ tới commit cha của nó, nên toàn bộ lịch sử của một nhánh chính là chuỗi commit duyệt ngược được từ con trỏ nhánh. Nhờ thiết kế này, việc tạo nhánh trong Git gần như tức thì và tốn rất ít dung lượng. Cách tổ chức này khuyến khích lập trình viên tạo nhánh riêng cho từng tính năng nhỏ thay vì dồn mọi thứ lên một nhánh duy nhất.

Con trỏ đặc biệt \verb|HEAD| cho biết nhánh nào đang được lấy ra làm việc (\textit{checked out}). Khi tạo commit mới, Git tự động dịch con trỏ của nhánh hiện tại tới commit vừa tạo, còn \verb|HEAD| vẫn bám theo nhánh đó \cite{atlassiangit}. Hiểu rõ quan hệ giữa \verb|HEAD|, nhánh và commit giúp sinh viên tránh nhầm lẫn khi chuyển nhánh hoặc hợp nhất code.

Với project ASPX, nhóm nên thống nhất quy ước đặt tên nhánh ngay từ đầu: nhánh tính năng bắt đầu bằng \verb|feature/|, nhánh sửa lỗi bắt đầu bằng \verb|fix/|, nhánh tài liệu bắt đầu bằng \verb|docs/|. Ví dụ, nhánh \verb|feature/login-aspx| chứa chức năng đăng nhập cho trang Web Forms, còn nhánh \verb|fix/webconfig-charset| sửa lỗi font chữ trong tệp cấu hình. Quy ước này giúp pull request ngắn gọn và pipeline kiểm thử chạy đúng phạm vi \cite{gitflow}.

\begin{lstlisting}[language=bash, caption={Tạo và quản lý nhánh tính năng}, label={lst:c23-branch}]
# List branches, star marks current branch
git branch
git status -sb
# Create and switch to feature branch (classic way)
git checkout -b feature/login-aspx
# Modern equivalent, needs Git 2.23 or newer
git switch -c feature/login-aspx
# Work, then commit on the branch
git add Default.aspx Default.aspx.cs
git commit -m "feat: add login form to Default page"
# Switch back to main and list all branches
git switch main
git branch -a
# Delete merged branch safely (-D forces delete)
git branch -d feature/login-aspx
\end{lstlisting}

Khối lệnh trên bao quát vòng đời điển hình của một nhánh: tạo nhánh bằng \verb|checkout -b| hoặc \verb|switch -c|, commit công việc lên nhánh, quay về nhánh \verb|main| rồi xóa nhánh đã hợp nhất. Tùy chọn \verb|-d| chỉ xóa khi nhánh đã được hợp nhất, trong khi \verb|-D| xóa cưỡng bức và có thể làm mất commit chưa merge. Bảng sau tóm tắt các lệnh hay dùng nhất khi làm việc với nhánh trong bài tập nhóm.

\begin{table}[H]
\centering
\caption{Các lệnh làm việc với nhánh}
\label{tab:c23-branch-cmd}
\begin{tabular}{ll}
\toprule
Lệnh & Ý nghĩa \\
\midrule
\texttt{git branch} & Liệt kê các nhánh local \\
\texttt{git branch -a} & Liệt kê cả nhánh local và remote \\
\texttt{git checkout -b ten-nhanh} & Tạo nhánh mới và chuyển sang ngay \\
\texttt{git switch ten-nhanh} & Chuyển sang nhánh đã có \\
\texttt{git switch -c ten-nhanh} & Tạo nhánh mới và chuyển sang ngay \\
\texttt{git branch -d ten-nhanh} & Xóa nhánh đã hợp nhất xong \\
\texttt{git branch -D ten-nhanh} & Xóa cưỡng bức nhánh chưa hợp nhất \\
\bottomrule
\end{tabular}
\end{table}

\begin{figure}[H]
\centering
\begin{tikzpicture}[>=Stealth, thick, box/.style={draw, rounded corners, fill=blue!10, minimum width=2.4cm, minimum height=0.9cm, align=center, font=\small}]
\node[box] (c1) {C1};
\node[box, right=1.2cm of c1] (c2) {C2};
\node[box, right=1.2cm of c2] (c3) {C3 \\ main};
\node[box, below=1.0cm of c2] (f1) {F1};
\node[box, right=1.2cm of f1] (f2) {F2 \\ login-aspx};
\node[above=0.2cm of c3, font=\footnotesize] {HEAD};
\draw[->] (c1) -- (c2);
\draw[->] (c2) -- (c3);
\draw[->] (c2) -- (f1);
\draw[->] (f1) -- node[midway, above, font=\footnotesize] {commit} (f2);
\draw[->, dashed] (f2) -- node[midway, right, font=\footnotesize] {merge} (c3);
\end{tikzpicture}
\caption{Mô hình hai nhánh và phép hợp nhất}
\label{fig:c23-branch-merge}
\end{figure}

Hình trên minh họa nhánh \verb|main| đi qua các commit C1, C2, C3 trong khi nhánh tính năng rẽ ra từ C2 với hai commit F1, F2 rồi hợp nhất trở lại. Mô hình rẽ nhánh và hợp nhất liên tục này là nền tảng của GitFlow và nhiều quy trình nhóm hiện đại \cite{gitflow}. Khi đã quen tay với nhánh, bước tiếp theo là hiểu sâu cơ chế hợp nhất và cách xử lý xung đột.

\subsection{Hợp nhất nhánh và giải quyết xung đột}

Lệnh \verb|git merge| hợp nhất lịch sử của một nhánh vào nhánh hiện tại, nhưng kết quả cụ thể phụ thuộc vào hình dạng lịch sử \cite{progit}. Nếu nhánh hiện tại chưa có commit mới nào kể từ điểm rẽ nhánh, Git chỉ cần trượt con trỏ tới trước, gọi là hợp nhất \textit{fast-forward}. Ngược lại, khi cả hai nhánh đều có commit mới, Git phải tạo một commit gộp (\textit{merge commit}) có hai commit cha. Trường hợp thứ hai gọi là hợp nhất ba bên (\textit{three-way merge}) và là dạng phổ biến nhất trong làm việc nhóm.

\begin{table}[H]
\centering
\caption{So sánh hợp nhất fast-forward và three-way}
\label{tab:c23-merge-compare}
\begin{tabular}{lll}
\toprule
Tiêu chí & Fast-forward & Three-way merge \\
\midrule
Điều kiện & Nhánh đích đứng yên từ điểm rẽ & Cả hai nhánh cùng có commit mới \\
Lịch sử & Tuyến tính, không commit gộp & Có commit gộp với hai cha \\
Lệnh thường dùng & \texttt{git merge ten-nhanh} & \texttt{git merge --no-ff ten-nhanh} \\
Hoàn tác & \texttt{git reset} về commit cũ & \texttt{git reset} hoặc \texttt{git revert} \\
Phù hợp & Nhánh cá nhân ngắn ngày & Nhánh tính năng của cả nhóm \\
\bottomrule
\end{tabular}
\end{table}

Xung đột (\textit{conflict}) xảy ra khi hai nhánh sửa cùng một vùng của cùng một tệp và Git không thể tự quyết định giữ bên nào. Để luyện tập, nhóm cố tình gây xung đột trên tệp \verb|Web.config| bằng cách cho hai nhánh sửa cùng khóa \verb|SiteName| trong khối \verb|appSettings|. Khối lệnh sau tái hiện toàn bộ quá trình ngay trên máy cá nhân của sinh viên.

\begin{lstlisting}[language=bash, caption={Gây xung đột có chủ đích trên tệp cấu hình}, label={lst:c23-merge-conflict}]
# Create feature branch and edit SiteName there
git switch -c feature/login-aspx
git commit -am "feat: rename site to shop on feature branch"
# Back on main, edit the SAME key differently
git switch main
git commit -am "feat: rename site to portal on main"
# Merge feature branch, Git stops with CONFLICT
git merge feature/login-aspx
# Inspect status and leftover whitespace errors
git status -sb
git diff --check
\end{lstlisting}

Sau lệnh hợp nhất, Git đánh dấu vùng xung đột trong \verb|Web.config| bằng ba loại dấu: \verb|<<<<<<< HEAD| mở đầu vùng của nhánh hiện tại, \verb|=======| ngăn cách hai bên, và \verb|>>>>>>> feature/login-aspx| kết thúc vùng của nhánh được hợp nhất. Không bao giờ commit nguyên các dấu này lên kho chung, vì tệp cấu hình sẽ hỏng và ứng dụng ASPX không khởi động được. Khối dưới đây cho thấy nội dung tệp tại thời điểm đang xung đột.

\begin{lstlisting}[language=XML, caption={Đoạn tệp cấu hình đang xung đột}, label={lst:c23-conflict-file}]
<?xml version="1.0" encoding="utf-8"?>
<configuration>
  <appSettings>
<<<<<<< HEAD
    <add key="SiteName" value="AspxDemo Portal" />
=======
    <add key="SiteName" value="AspxDemo Shop" />
>>>>>>> feature/login-aspx
  </appSettings>
</configuration>
\end{lstlisting}

Các bước giải quyết xung đột chuẩn cho sinh viên gồm:
\begin{enumerate}
\item Mở \verb|Web.config|, đọc kỹ cả hai phía của vùng xung đột để hiểu ý đồ của mỗi bên.
\item Giữ lại một phía hoặc viết lại nội dung gộp cả hai ý, rồi xóa sạch ba loại dấu xung đột.
\item Chạy \verb|git diff --check| để chắc chắn không còn dấu thừa hay khoảng trắng lỗi.
\item Đánh dấu đã giải quyết bằng \verb|git add Web.config| rồi hoàn tất với \verb|git commit|.
\item Kiểm tra lại bằng \verb|git status| và \verb|git log --graph| trước khi đẩy lên kho chung.
\end{enumerate}

\begin{lstlisting}[language=XML, caption={Tệp cấu hình sau khi giải quyết xung đột}, label={lst:c23-resolved-file}]
<?xml version="1.0" encoding="utf-8"?>
<configuration>
  <appSettings>
    <add key="SiteName" value="AspxDemo Portal" />
  </appSettings>
</configuration>
\end{lstlisting}

Nếu đang hợp nhất dở mà nhận ra mình thao tác nhầm nhánh, lệnh \verb|git merge --abort| sẽ hủy toàn bộ tiến trình và đưa cây làm việc về trạng thái trước khi merge. Sau khi hủy, nên chạy \verb|git status| để xác nhận mọi thứ sạch sẽ rồi mới thao tác lại \cite{atlassiangit}. Với xung đột phức tạp, sinh viên có thể nhờ công cụ gộp bằng hình ảnh của Visual Studio hỗ trợ, nhưng vẫn phải đọc hiểu từng vùng trước khi bấm chấp nhận.

\subsection{Tệp .gitignore chuẩn cho dự án ASP.NET}

Tệp \verb|.gitignore| liệt kê các đường dẫn mà Git phải bỏ qua, giúp kho chứa mã nguồn gọn nhẹ và tránh rò rỉ tệp của máy cục bộ \cite{githubdocs}. Với project ASP.NET Web Forms, ít nhất năm nhóm phải luôn bị bỏ qua: thư mục biên dịch \verb|bin/| và \verb|obj/|, thư mục gói \verb|packages/|, thư mục cấu hình Visual Studio \verb|.vs/|, và tệp tùy chọn người dùng \verb|*.user|. Khối mẫu đầy đủ dưới đây có thể sao chép nguyên văn vào mọi đồ án ASPX của môn học. Mỗi khối đều có chú thích để sinh viên hiểu rõ tác dụng trước khi áp dụng.

\begin{lstlisting}[language=bash, caption={Mẫu gitignore đầy đủ cho ASP.NET Web Forms}, label={lst:c23-gitignore}]
# Build outputs, rebuilt by Visual Studio or MSBuild
bin/
obj/

# Visual Studio per-machine files
.vs/
*.user
*.suo

# NuGet packages, restored on build
packages/
*.nupkg

# Publish profiles may leak server paths
*.pubxml
PublishProfiles/

# OS junk and logs
Thumbs.db
.DS_Store
*.log
*.tmp
\end{lstlisting}

\begin{description}
\item[Build] Hai mục \verb|bin/| và \verb|obj/| loại bỏ toàn bộ kết quả biên dịch, vì chúng được tạo lại sau mỗi lần build.
\item[Visual Studio] Nhóm \verb|.vs/|, \verb|*.user| và \verb|*.suo| chứa trạng thái cửa sổ của từng máy, không có giá trị chia sẻ cho cả nhóm.
\item[NuGet] Thư mục \verb|packages/| và \verb|*.nupkg| rất nặng nên chỉ lưu danh sách gói, khi build hệ thống sẽ khôi phục tự động.
\item[Publish] Hồ sơ phát hành có thể lộ đường dẫn máy chủ, vì vậy \verb|*.pubxml| và \verb|PublishProfiles/| cần bị chặn triệt để.
\item[Hệ điều hành] Các tệp \verb|Thumbs.db|, \verb|.DS_Store| và \verb|*.log| do hệ điều hành và công cụ tự sinh, không liên quan tới mã nguồn.
\end{description}

Lưu ý quan trọng: \verb|.gitignore| chỉ có tác dụng với tệp chưa được theo dõi. Nếu ai đó đã commit nhầm \verb|bin/| lên kho chung, việc thêm quy tắc bỏ qua sau đó không xóa được lịch sử cũ, và đây cũng là câu hỏi phản biện quen thuộc của giảng viên. Cách sửa là gỡ tệp khỏi chỉ mục nhưng giữ lại trên đĩa bằng \verb|git rm --cached|, rồi commit kèm \verb|.gitignore| mới như khối lệnh dưới đây.

\begin{lstlisting}[language=bash, caption={Gỡ tệp đã commit nhầm ra khỏi chỉ mục}, label={lst:c23-rm-cached}]
# Check which rule ignores a path, if any
git check-ignore -v Default.aspx
# File was committed before gitignore existed
git ls-files | findstr "bin obj user"
# Untrack it, but keep the local copy on disk
git rm --cached -r bin obj
git rm --cached AspxDemo.csproj.user
# Commit the fix together with gitignore
git add .gitignore
git commit -m "fix: stop tracking build outputs and user files"
git status -sb
\end{lstlisting}

\subsection{Dự án mini với ASPX: tạo trang, commit và xem lịch sử}

Để gắn kiến thức nhánh và hợp nhất với học phần, nhóm xây dựng một project mini gồm đúng một trang \verb|Default.aspx| và tệp code-behind \verb|Default.aspx.cs| trên .NET Framework 4.8.1 \cite{msnetfx}. Trang hiển thị một nhãn chào mừng và một nút bấm, đủ để demo vòng đời chỉnh sửa, commit, rẽ nhánh và hợp nhất. Mọi chuỗi chữ trong ví dụ đều viết không dấu để giữ mã nguồn hoàn toàn ASCII.

\begin{lstlisting}[language=XML, caption={Trang Default ASPX cho dự án mini}, label={lst:c23-default-aspx}]
<%@ Page Title="Trang chu" Language="C#" AutoEventWireup="true" CodeBehind="Default.aspx.cs" Inherits="AspxDemo._Default" %>
<!DOCTYPE html>
<html xmlns="http://www.w3.org/1999/xhtml">
<head runat="server">
    <title>AspxDemo Mini</title>
</head>
<body>
    <form id="form1" runat="server">
        <div>
            <asp:Label ID="lblMessage" runat="server" Text="Xin chao ASP.NET" />
            <br />
            <asp:Button ID="btnClick" runat="server" Text="Nhan vao day" OnClick="btnClick_Click" />
        </div>
    </form>
</body>
</html>
\end{lstlisting}

Chỉ thị \verb|Page| ở dòng đầu khai báo ngôn ngữ C\# và liên kết tới tệp code-behind \verb|Default.aspx.cs| qua thuộc tính \verb|CodeBehind|. Hai điều khiển \verb|Label| và \verb|Button| đều có \verb|runat="server"| nên truy cập được từ code-behind. Khi bấm nút trên trình duyệt, sự kiện \verb|btnClick_Click| sẽ chạy trên máy chủ rồi trả HTML mới về cho người dùng.

\begin{lstlisting}[language=[Sharp]C, caption={Tệp code-behind xử lý sự kiện nút bấm}, label={lst:c23-codebehind}]
using System;
namespace AspxDemo
{
    public partial class _Default : System.Web.UI.Page
    {
        protected void Page_Load(object sender, EventArgs e)
        {
            if (!IsPostBack)
            {
                // First visit, show default greeting
                lblMessage.Text = "Xin chao ASP.NET 4.8.1";
            }
        }
        protected void btnClick_Click(object sender, EventArgs e)
        {
            // Button click, update the label text
            lblMessage.Text = "Ban vua nhan nut, chuc mung!";
        }
    }
}
\end{lstlisting}

Lớp \verb|_Default| kế thừa từ trang Web Forms và xử lý hai sự kiện: \verb|Page_Load| hiển thị lời chào khi mở trang lần đầu, còn \verb|btnClick_Click| đổi nội dung nhãn khi người dùng bấm nút. Thuộc tính \verb|IsPostBack| giúp phân biệt lần mở trang đầu tiên với các lần gửi biểu mẫu về máy chủ. Sau khi lưu hai tệp, sinh viên chạy khối lệnh Git dưới đây để commit rồi quan sát lịch sử phân nhánh.

\begin{lstlisting}[language=bash, caption={Commit dự án mini và xem lịch sử phân nhánh}, label={lst:c23-demo-log}]
# Stage the new page and commit on main
git add Default.aspx Default.aspx.cs Web.config
git commit -m "feat: create Default page with greeting button"
# Create branch, edit button text, commit there
git switch -c feature/login-aspx
git commit -am "feat: change button text on login branch"
# Merge back to main and view the graph
git switch main
git merge --no-ff feature/login-aspx -m "Merge branch 'feature/login-aspx'"
git log --graph --oneline --decorate -n 6
# Sample output of the log command above
# * 9f3a1c2 (HEAD -> main) Merge branch 'feature/login-aspx'
# |\
# | * 7b2e4a1 feat: change button text on login branch
# | * 2c1d8f0 feat: create Default page with greeting button
# |/
\end{lstlisting}

Tùy chọn \verb|--graph| vẽ sơ đồ phân nhánh bằng ký tự ASCII, còn \verb|--oneline| và \verb|--decorate| giúp mỗi commit gọn trong một dòng kèm tên nhánh. Thông điệp commit tuân thủ Conventional Commits với tiền tố \verb|feat:| và \verb|docs:| giúp lịch sử dễ đọc và dễ sinh ghi chú phát hành \cite{convcommits}. Hoàn thành dự án mini này, sinh viên đã đủ kỹ năng chuyển sang các thao tác nâng cao như hoàn tác, stash và gắn thẻ phiên bản.